% =========================================================================
%  The Normativity of Self-Legislation in Intimate Relationships:
%  Freedom, Morality, and Justice
%  Philosophy of Intimacy and the Theory of Justice — Paper III
%  Wanhong Huang (Independent researcher)
%  Compile: xelatex main.tex && biber main && xelatex main.tex && xelatex main.tex
% =========================================================================
\documentclass[11pt,a4paper]{article}
\usepackage{serendip-paper}

% ── Bibliography (add refs.bib to this folder) ───────────────────────
\usepackage[backend=biber, style=authoryear, sorting=nyt, maxcitenames=2]{biblatex}
\addbibresource{refs.bib}

\begin{document}

% ── Title page ───────────────────────────────────────────────────────
\thispagestyle{empty}
\null\vfill
\begin{center}
{\footnotesize\color{rosedeep}\scshape Philosophy of Intimacy and the Theory of Justice \quad$\cdot$\quad Paper III\par}
\vspace{0.9em}
\coverrule\par
\vspace{1.6em}
{\color{warnred}\Large The Normativity of Self-Legislation in Intimate Relationships: Freedom, Morality, and Justice\par}
\vspace{1.0em}
{\color{warnred}\large\itshape On a Lover's Vow, the Rights-Order It Institutes, and the Conditions under Which It Becomes Domination\par}
\vspace{2.4em}
\coverrule[0.25]\par
\vspace{1.4em}
{\large Wanhong Huang\par}
\vspace{0.5em}
{\footnotesize \texttt{huangwanhong.g.official@gmail.com}\par}
\vspace{1.2em}
{\today\par}
\end{center}
\vfill\null
\coverfootnote
\clearpage

% ── Dedication ───────────────────────────────────────────────────────
\clearpage
\thispagestyle{empty}
\vspace*{0.20\textheight}
\begin{center}
{\itshape\color{ink}
To the one I love,\\[0.4em]
who taught me that to keep a vow\\[0.25em]
is not to hold a course against her,\\[0.25em]
but to keep, between us, the room in which she may always answer.\\}
\vspace{2.6em}
{\cjkfont\large\color{rosedeep} 执子之手，与子偕老}\\[0.8em]
{\itshape\small\color{rosedeep}
I take your hand,\\[0.18em]
and with you I will grow old.}\\[1.0em]
{\footnotesize\color{lightgold}from the \cjkfont 诗经\/{\normalfont\ (\textit{Book of Songs})}}
\vspace{2.4em}

{\cjkfont\large\color{rosedeep} 发乎情，止乎礼义}\\[0.8em]
{\itshape\small\color{rosedeep}
It rises from feeling,\\[0.18em]
and comes to rest in what is right.}\\[1.0em]
{\footnotesize\color{lightgold}after the \cjkfont 毛诗序\/{\normalfont\ (Great Preface to the \textit{Book of Songs})}}
\end{center}
\vfill
\clearpage

% =========================================================================
\begin{abstract}
\noindent
Consider a lover who legislates for himself an unconditional law: \emph{whatever happens, however the beloved feels, I will keep loving her.} This paper subjects such a vow to a three-layered normative analysis, and uses it as a lens on a general structure: the conditions under which the self-legislation of a shared life is free, the conditions under which it is just, and the conditions under which it silently becomes domination or exploitation.

\smallskip
\noindent The argument is integrative. (1) On \emph{freedom}: the vow is exemplary Kantian autonomy, but Kant's own apparatus, the fact of reason and the noumenal self, cannot tell us \emph{when} a self-legislation is free; a compatibilist, reasons-responsive criterion can, and on it the vow's freedom turns out to be conditional on its remaining responsive to the beloved as a source of reasons. (2) On \emph{morality}: the vow is admirable as constancy, yet its content, what \emph{loving her} means, is left to the lover alone to fix, in tension with the responsiveness that the philosophy of love takes to be constitutive of love. (3) On \emph{justice}: drawing on republican non-domination and a Hohfeldian analysis of the rights-order a vow institutes, the paper argues that a unilaterally legislated law over a two-person life cannot be just, because justice between two agents is modally, not merely actually, a matter of each retaining the standing to contest the terms.

\smallskip
\noindent A dedicated section on the constitution of the legislating subject (structuralist and feminist) shows why the injustice is not idiosyncratic but structurally produced: the ``self'' that legislates is itself constituted by a historically gendered order, so that a sincere, universalizable vow can reproduce that order's subordinations. Integrating Marxian and social-reproduction theory, the paper distinguishes the conditions under which the vow institutes a legitimate rights-order from those under which it institutes \emph{exploitation}, the appropriation of another's unrecognized relational and reproductive labour under an ideology that presents the appropriation as love. Against the objection that turning to relationality merely relocates the problem, since recognition can itself be ideological, it argues that the operative criterion is the \emph{live contestability} of the terms, not their de facto endorsement. The conclusion is an isomorphism: the micro-order of a vow and the macro-order of a political economy share a structure, and in both the standing of the legislation rests not on the purity of the legislator but on the irreducible, modally robust standing of the other to answer.

\vspace{0.6em}
\noindent\textbf{Keywords}\quad
self-legislation; autonomy; reasons-responsiveness; non-domination; recognition; ideological recognition; Hohfeldian rights; social reproduction; exploitation; relational autonomy; justice in intimate relationships.
\end{abstract}

\tableofcontents
\clearpage

% =========================================================================
\section{Introduction: A Vow and Its Puzzle}\label{sec:intro}

Consider a lover who, reflecting on the precariousness of feeling and the contingency of circumstance, resolves upon a law for himself:

\begin{quote}\itshape
Whatever happens, however she feels toward me, I will keep loving her.
\end{quote}

The resolution is not a prediction about how he expects to feel. It is a self-imposed principle, a law the agent gives himself and intends to govern his conduct regardless of what the world, or the beloved, does. It is, in the most literal sense, an act of \emph{self-legislation}: the will binding itself to a maxim it has authored. Such vows are among the most admired achievements of a loving life. And yet, examined closely, the vow harbours a puzzle, and the puzzle, once unfolded, turns out to be a window onto a structure far larger than itself.

On the one hand, the vow is a triumph of \emph{freedom}: it liberates the lover's love from bondage to the beloved's moods, to reciprocation, to the favourable alignment of fortune. On the other hand, the very structure that wins this freedom, the love's unconditionality, its insulation from how the beloved actually responds, raises a worry not about freedom at all, but about \emph{justice}. A relationship is not a solitary undertaking; it has two parties; and a law that one party legislates unilaterally, however lovingly, governs a shared life the other has had no hand in authoring. Stated at its sharpest: the vow institutes, between two people, an order, an arrangement of who may expect what, who owes what, who may alter the terms. The question this paper pursues is when such an order, unilaterally legislated, is \emph{free}, when it is \emph{just}, and when it silently becomes something else, an order of domination, even of exploitation, wearing the face of love.

\subsection{The Thesis and the Isomorphism}

The thesis has a layered form. (i) The vow is exemplary self-legislation and, in the Kantian sense, an exercise of freedom; but Kantian freedom cannot say \emph{when} a self-legislation is free, and a more operable, compatibilist criterion reveals the vow's freedom to be conditional. (ii) The vow can be morally admirable, but its admirable form is silent on its content, and the unilateral fixing of that content is already in tension with what love is. (iii) The vow cannot, in its unilateral form, be \emph{just}, because justice between two parties is constitutively a matter of each retaining a modally robust standing to contest the terms, which a unilateral law by construction withholds. (iv) The injustice is not the failing of a bad lover but a structural product: the legislating ``self'' is itself constituted by a historically gendered order, so that even a sincere and universalizable vow can reproduce subordination. (v) There is a determinate line between the vow that institutes a legitimate rights-order and the vow that institutes \emph{exploitation}; the line is drawn by whether the other retains co-authorship over the terms of what is given and demanded, or is instead enclosed, under the ideology of love, as the unrecognized provider of relational labour. (vi) The remedy is not a purer legislator but a relational re-grounding in which the other's standing to contest is constitutive; and the criterion that keeps this from collapsing back into ideology is live contestability rather than endorsement.

The deepest claim is an \emph{isomorphism}. The micro-order a lover institutes by a vow and the macro-order a political economy institutes by its distribution of rights and labour share a structure. In both, a legislation that looks self-grounded is in fact grounded in relations that precede it; in both, procedural correctness on the legislating side cannot certify justice; in both, what distinguishes order from domination is whether those bound by the law retain a real, not merely nominal, power to contest it. The intimate case is the clearest lens on this structure because its two-person character is vivid and because the lover's sincerity strips away any suspicion of mere bad faith: if even \emph{this} can be domination, the diagnosis cannot be dismissed.

\subsection{Method and Scope}

The vow is treated as a first-person practical datum, in the manner of \textcite{frankfurt2004reasons} and \textcite{velleman1999love}, rather than as a third-personal object. The lover is assumed sincere, reflective, and well-meaning throughout; the argument is most secure precisely in that best case. The paper integrates traditions that rarely appear together, Kantian autonomy, compatibilist free-will theory, republican non-domination, analytic jurisprudence, structuralist and feminist accounts of subject-formation, and Marxian social-reproduction theory, but it does so under the discipline of a single spine, stated above, so that the integration is argumentative rather than encyclopaedic. Two limitations are flagged at the outset. First, the paper takes from Marx the \emph{formal structure} of exploitation and the thesis of ideological masking, while remaining agnostic about the labour theory of value; nothing in the argument requires that theory to be true. Second, the constructive criterion offered at the end, live contestability, is defended as the right criterion, not as an algorithm; its residual hard cases are acknowledged rather than concealed.

The plan follows the spine. Section~\ref{sec:kant} reconstructs Kantian self-legislation and exposes its internal difficulty about the grounds of obligation. Section~\ref{sec:freewill} asks when a self-legislation is \emph{free} and develops a reasons-responsive criterion. Section~\ref{sec:moral} treats the moral layer and the first crack. Section~\ref{sec:justice} develops the justice layer through non-domination. Section~\ref{sec:rights} analyses the vow as the institution of a Hohfeldian rights-order. Section~\ref{sec:subject} examines the constitution of the legislating subject. Section~\ref{sec:exploitation} draws the line between rights-order and exploitation. Section~\ref{sec:ideology} meets the objection from ideological recognition. Section~\ref{sec:regrounded} reconstructs the vow. Section~\ref{sec:conclusion} states the isomorphism.

% =========================================================================
\section{Kantian Self-Legislation and the Problem of Its Authority}\label{sec:kant}

The vow's claim to be an exercise of freedom rests on the Kantian distinction between autonomy and heteronomy. A will is \emph{heteronomous} when determined by something external to its own rational legislation, by inclination, desire, circumstance, or the anticipated favour of others; it is \emph{autonomous} when the law it follows issues from itself, when the will is ``a law to itself'' \parencite[4:440]{kant1785groundwork}. Autonomy is not the absence of law but the presence of \emph{self-given} law; freedom is not doing whatever one is moved to do, which is precisely the condition of being moved by alien causes, but acting on principles one has authored as a rational agent.

This yields the apparent paradox that self-binding is freedom. The agent who acts on whichever impulse is strongest is not free but driven; the agent who binds himself to a self-authored principle achieves a standing independence from the flux of impulse and becomes the author of his conduct rather than its venue. \textcite{korsgaard1996sources} reconstructs this as the claim that to act autonomously is to act on a law one gives oneself in light of a \emph{practical identity}, and that this self-legislation is what constitutes one as a unified agent rather than a succession of psychological states; the thought is extended in \textcite{korsgaard2009self}, where agency itself is an achievement of self-constitution under self-given principles.

Applied to the vow: the lover who loves only while loved in return, only while moods are favourable, only while circumstance cooperates, has a love that is heteronomously conditioned, the dependent variable of factors outside his legislation. The vow severs these dependencies. In resolving to love ``whatever happens, however she feels,'' the lover withdraws his love from the jurisdiction of mood and reciprocity and places it under his own legislation. On Kantian grounds the vow is therefore not a diminishment of freedom but its fullest expression, and there is a recognizable Kantian thought that conduct from principle has a higher worth than conduct from inclination, the former being the agent's achievement, the latter a gift of temperament \parencite[4:398--399]{kant1785groundwork}.

\subsection{The Fact of Reason and the Circle of Authority}

Grant all this, and a question remains that will prove decisive: \emph{whence the authority of the self-given law?} If the lover is both the author of the law and the one bound by it, what prevents him from repealing it the moment it chafes? A law one can revise at will is no law; yet a law one gives oneself seems to be exactly that. Kant felt the difficulty acutely. In the \emph{Groundwork} he sought to derive the moral law's bindingness from the very concept of a rational will, but by the \emph{Critique of Practical Reason} he rests the matter on what he calls the \emph{fact of reason} (\emph{Faktum der Vernunft}): our consciousness of the moral law as binding is a fact that reason simply presents, neither derived from anything prior nor in need of, nor susceptible of, external proof \parencite[5:31]{kant1788practical}. The bindingness of self-legislation is thus secured by treating it as a datum rather than a conclusion.

This is the first of two pressure points. To ground the authority of self-given law in a ``fact'' that reason presents to itself is either to assert what was to be shown, that the self-given law really does bind its author, or to relocate the authority outside the individual will, into ``reason'' as something the agent does not control and cannot repeal. Kant takes the second route: the law binds not because \emph{this} individual wills it but because reason, in him, legislates it, and reason is not his to revoke. The move saves the bindingness but at a price we will collect repeatedly: it makes the genuinely authoritative legislator not the empirical person but \emph{reason as such}, and thereby quietly concedes that the individual, considered as the contingent being he is, is not in fact the self-sufficient source of the law he appears to give himself.

\subsection{The Noumenal Refuge}

The second pressure point is the home of Kant's answer to every charge that the legislator is corruptible. Confronted with the worry that an agent's sense of what reason requires might be distorted by inclination, history, or circumstance, Kant can reply that such distortions belong to the \emph{empirical} self, the self in the world of appearances, whereas the true legislator is the \emph{noumenal} self, the rational agent considered as a member of the intelligible order, whose legislation is by its nature untainted by empirical determination. Freedom, on this picture, is a property of the noumenal will; the moral law is what that will gives itself; and the contingencies that might corrupt an actual person's judgement are, in principle, screened off from the legislating self that matters.

This noumenal refuge is the strongest move available to the defender of pure self-legislation, and much of this paper can be read as the attempt to show that it cannot do the work the intimate case demands of it. Two observations begin that work. First, the refuge is purchased by placing the authoritative legislator beyond the reach of experience, which makes it equally beyond the reach of any \emph{operable} test: if the question is \emph{when}, in actual life, a given self-legislation is free rather than driven, just rather than dominating, an appeal to a noumenal self that is by definition never given in experience yields no usable answer. Second, the refuge concedes, in the very act of invoking it, that the \emph{empirical} legislator, the actual lover making the actual vow, is exposed to exactly the distortions the paper will track; it protects the purity of self-legislation only by relocating it to a self that does no concrete work. We will let the refuge stand for now and return, in Section~\ref{sec:subject}, to dismantle the assumption it most depends on: that there is a legislating self constituted prior to, and independently of, the relations and history in which the empirical person is formed.

\subsection{Korsgaard's Half-Step}

Korsgaard's reconstruction is important here precisely because it takes a half-step away from the noumenal refuge and toward the relational view, without completing it. By grounding obligation in \emph{practical identities}, the descriptions under which we value ourselves and find our lives worth living, she brings the source of normativity down from the noumenal into the texture of an agent's actual self-understanding \parencite{korsgaard1996sources}. Practical identities are largely social: one is a friend, a citizen, a parent, a lover, and these identities, with their attendant obligations, are not legislated \emph{ex nihilo} by a noumenal will but inherited, adopted, and sustained in a social world. To that extent Korsgaard has already conceded that the legislator is socially constituted. Yet she retains a backstop: beneath all contingent practical identities lies one that is not optional, our identity simply as reflective rational agents, members of the ``kingdom of ends,'' and it is this that supplies universal moral obligations and rescues the account from relativism. The half-step is therefore arrested: the social constitution of the self is admitted at every level but the last, where a non-relational, universal identity is reinstated to do the grounding work. The argument of this paper presses on exactly that last level. If even the deepest practical identity, the sense of oneself as one who loves, and of what loving requires, is constituted in and by relations that are historically structured, then there is no non-relational floor on which a self-sufficient legislation could stand. Section~\ref{sec:subject} makes that case; the intervening sections show why it matters.

% =========================================================================
\section{When Is a Self-Legislation Free? From the Noumenal to the Reasons-Responsive}\label{sec:freewill}

Section~\ref{sec:kant} establishes that the vow is, in Kantian terms, an exercise of autonomy. But it also exposes that Kantian autonomy cannot answer the question the rest of the paper needs answered: \emph{when} is a particular act of self-legislation free, and when is it merely a sophisticated form of being driven? The noumenal account locates freedom in a self never given in experience and so offers no operable criterion; the fact of reason secures bindingness by fiat. To make progress we need a conception of freedom that applies to the legislating mechanism as it actually operates. The compatibilist tradition in the theory of free will supplies one.

\subsection{Two Compatibilist Resources}

Two resources are pertinent. The first is Frankfurt's hierarchical account, on which an agent is free in the relevant sense when his will, the desire that moves him to act, is one he reflectively endorses at a higher order: the unwilling addict, moved by a first-order desire he repudiates at the second order, is unfree, while the agent whose effective desire is ratified by his higher-order volitions acts of his own free will \parencite{frankfurt1971freedom}. The second, more demanding and more useful here, is Fischer and Ravizza's account of \emph{guidance control}, which they analyse as two conditions: the action must issue from the agent's \emph{own} mechanism (ownership), and that mechanism must be \emph{moderately reasons-responsive}, that is, disposed to recognize and respond to a suitable range of reasons, including reasons to do otherwise, across a range of scenarios \parencite{fischer1998responsibility}. Notably, their theory is \emph{historicist}: whether a mechanism is the agent's own depends on how it came to be his, on the history of its formation, so that a mechanism implanted by manipulation, however well it functions now, does not ground freedom. We will exploit this historicism in Section~\ref{sec:subject}.

\subsection{Applying Reasons-Responsiveness to the Vow}

On the reasons-responsive criterion, the freedom of the vow is not settled by the bare fact that the lover gave himself the law. It depends on whether the mechanism that now governs his loving remains responsive to the relevant reasons. And here the central reasons are supplied by the beloved: her needs, her changes, her objections, her saying that what he offers as love is not, for her, love. A vow whose whole point is to hold steady ``however she feels'' is, read in one way, a mechanism engineered to be \emph{un}responsive to precisely these reasons, to treat her dissent not as input that could redirect the loving but as weather to be outlasted. To that extent the vow, far from being the height of freedom, approaches the condition of the driven: the lover is moved by a standing resolution that has been insulated against the reasons that ought to be able to move it. The Kantian sees only that the law is self-given and pronounces the vow free; the reasons-responsive criterion sees that a self-given law can be self-given as a cage, and asks the further question the Kantian cannot.

This yields a determinate and somewhat surprising result. The vow is free \emph{to the extent that}, and \emph{only to the extent that}, the constancy it institutes leaves the lover's loving responsive to the beloved as a source of reasons. A constancy of \emph{presence and commitment} that remains responsive in its \emph{mode} is free; a constancy that fixes the mode and immunizes it against her is, by the very criterion that makes self-legislation a candidate for freedom, a diminution of freedom dressed as its triumph. Freedom and the seed of the later injustice are thus located at a single point: the responsiveness of the legislating mechanism to the other.

\subsection{The Hole Both Theories Share}

Reasons-responsiveness is a real advance, but it has a famous limitation that, far from being a defect for our purposes, is the hinge of the whole integration. The criterion is officially neutral about \emph{which} reasons the mechanism must answer to and about \emph{who} has standing to supply them; it asks only that the mechanism be suitably responsive to whatever reasons there are. As critics have noted, this leaves the account silent on the \emph{sources} of the reasons. An agent may be exquisitely responsive to the reasons he recognizes while being systematically closed to a whole class of reasons he does not recognize as reasons at all, for instance, the claims of someone whose standing to make claims he has, without noticing, discounted in advance. Such an agent passes the reasons-responsiveness test on the reasons that reach him, yet is, in a deeper sense, unfree, and, as we will see, positioned to dominate.

This is the same hole, arrived at from the other side, that Kantian universalizability was meant to plug and could not. Kant's test asks whether one could will one's maxim as a universal law; the hope was that this formal demand would force the agent to take account of everyone's standing. Hegel's objection, that the test screens out only outright contradictions and licenses a wide range of substantively objectionable maxims provided they are described with care \parencite[§135]{hegel1820right}, is precisely the complaint that universalizability does not in fact secure attention to the other's standing: a maxim that discounts another's claims in advance can be universalized without contradiction, because the discounting has already been built into the description of the situation. So the Kantian formal route and the compatibilist reasons-responsive route fail at the \emph{same} place: neither can, from its own resources, guarantee that the legislating self is open to the other as a full source of reasons and claims. Both presuppose a legislator already disposed to count the other, and neither can supply that disposition. That common gap is where relationality must enter, not as a preference but as the only thing that can fill a hole both leading theories of free and responsible agency leave open. The remainder of the paper is, in effect, the filling of that hole.

% =========================================================================
\section{The Moral Layer and the First Crack}\label{sec:moral}

Does the vow's standing as freedom carry into moral admirability? On several theories, initially yes. On a Kantian reckoning it honours the Formula of Humanity: the lover commits to treating the beloved never merely as a means, to be retained while pleasing and discarded when not, but as an end whose standing does not fluctuate with her usefulness \parencite[4:429]{kant1785groundwork}. There is, too, a venerable view that unconditional commitment is close to love's essence. On Singer's bestowal account, love confers a value not earned by the beloved's merits and so not revoked when they waver \parencite{singer1991nature}. On Frankfurt's account, love is a mode of volitional necessity in which the lover comes to be unable to be indifferent to the beloved's flourishing, and this bondage is the very shape of caring \parencite{frankfurt2004reasons}. By these lights the vow articulates love rather than distorting it.

Yet here the first crack appears. The vow's content is ``I will keep loving her'', but this is schematic until we specify what \emph{loving her} consists in, and the unilateral structure of self-legislation lets the lover, and the lover alone, fill in that content. ``Keep loving her'' can mean: \emph{I will keep willing her good, attending to her as the particular person she is, responding to what she actually needs and asks.} Or it can slide, by degrees, into: \emph{I will keep doing what I have determined loving her to be, whatever she says about it.} The first is responsive; the second, despite its devotion, is imposition.

The philosophy of love supplies the resources to see why the second reading is defective \emph{as love}. For Velleman, love is a moral emotion, the lover's response to the beloved as a rational agent, an arresting awareness of her worth as a being with her own will and ends \parencite{velleman1999love}. A love whose content is fixed unilaterally and held immune to the beloved's actual will is in tension with this: it addresses not the beloved as a rational agent with claims of her own, but an image of her, or the lover's conception of her good, around which the conduct of love is then organized. Kolodny's relational theory presses the same point: love is grounded in, and answerable to, an ongoing relationship between two people, not in a disposition lodged in one and sealed against the other \parencite{kolodny2003love}; and Helm's account of intimate identification likewise makes the beloved's own perspective constitutive of the loving, not an optional input to it \parencite{helm2010love}. Even Frankfurt's volitional necessity, properly understood, is a necessity to care about \emph{the beloved's} good as she has it, not about the lover's project of loving.

So the crack is this. Stated abstractly, the vow is admirable; but its abstract statement conceals a fork, and nothing in the \emph{form} of self-legislation determines which tine is taken, because the form is silent on content. Read responsively, the vow remains love and remains admirable. Read as the unilateral fixing of love's content, it converts the beloved from the addressee of love into the occasion for the lover's exercise of a self-assigned role. The question the freedom and morality layers cannot answer is now unavoidable, and it is a question about \emph{authority}: who decides what loving her means? That is a question of justice.

% =========================================================================
\section{The Justice Layer: Non-Domination and the Limits of Unilateral Legislation}\label{sec:justice}

The deepest assessment concerns neither the vow's freedom nor its moral quality from the lover's side, but its \emph{justice}, the standing of the arrangement it institutes between two people. I begin from a deliberately thin and widely shared conception and let it do most of the work, reserving a thicker conception for Section~\ref{sec:regrounded}.

\subsection{Why Procedure Cannot Certify Content}

In the moral philosophy descended from Kant there is a standing question whether a correct \emph{procedure} of legislation guarantees just \emph{content}. Kant's wager is that it does: a maxim that survives universalizability, willed by a rational agent treating humanity as an end, is thereby in order, so that one need not consult an independent standard of justice to certify the result. The wager has been contested from two directions. Hegel's charge of empty formalism, already invoked, is that universalizability screens out only contradictions and permits much that is substantively unjust \parencite[§135]{hegel1820right}. Rawls, inheriting the Kantian ambition, declined to trust the bare form and instead built a procedure, the original position behind a veil of ignorance, whose design \emph{embeds} substantive constraints, equality, mutual disinterest, ignorance of one's place, precisely so that its outputs would be just \parencite{rawls1971theory}. The lesson is double-edged: a procedure can be made to guarantee just content, but only by loading morally significant conditions into it in advance. Procedure does not generate justice for free.

A further consideration is decisive for our case. Even a procedurally impeccable self-legislation can produce unjust content when the legislator's own judgement, of what love requires, what a partner owes, what counts as a reasonable claim by the other, has been formed by conditions that systematically under-weight the other's standing. The procedure then runs correctly on corrupted inputs. In the intimate sphere this is not a remote possibility but the ordinary condition, since everyone's conception of love is inherited and acculturated; Section~\ref{sec:subject} shows how. For now the structural point suffices: the lover's sincerity, freedom, and good will establish that the vow is, from his side, autonomously and even admirably made, but they cannot establish that it is \emph{just}, because justice between two parties is not the kind of thing one party can confer by legislating well.

\subsection{Non-Domination as the Operative Conception}

The thin conception I adopt is republican: a relationship is unjust to the extent that one party is subject to \emph{domination}, understood, after Pettit, as exposure to another's capacity for arbitrary interference, interference not controlled by, and not answerable to, the affected party's own avowable interests and voice \parencite{pettit1997republicanism}. The decisive feature of this conception, for us, is that domination is \emph{modal}: one is dominated not only when interference actually occurs but whenever one is exposed to its possibility at another's discretion. Pettit's own illustration is the slave of a kindly master: the master who chooses never to interfere still dominates, because the non-interference depends on his continuing goodwill rather than on the slave's secured standing; the slave lives at discretion, and that is unfreedom even when the discretion is exercised benignly \parencite{pettit1997republicanism}. What makes the difference between freedom and domination is therefore not the quality of the treatment but the \emph{structure of the power}: whether the dominated party has a secured, non-discretionary standing, or merely the good fortune of a benevolent superior.

This is exactly the tool the intimate case requires, because the lover of our vow is, at his best, the benevolent master. His love may be impeccable in its actual exercise; the question non-domination teaches us to ask is whether the beloved's good treatment is \emph{secured} by her standing or merely \emph{granted} by his disposition. A unilaterally legislated law of love, however benevolent, places the beloved in the structural position of one who lives at the lover's discretion: the terms of the love, what it will demand of her, what it will provide, what it will refuse to hear, are set by him, and her good treatment depends on his continuing to will it rather than on any standing of hers to contest it. That is domination in the modal sense even if the lover never once interferes badly, indeed even if he is, by every actual measure, wonderful. The benevolent-master structure is the precise diagnosis of what is wrong with a vow that is, by hypothesis, made in love.

\subsection{Three Vectors of Domination}

The general diagnosis resolves into three concrete vectors.

\textbf{(i) Immunization against voice.} The vow's signature clause is ``however she feels.'' But the beloved's feelings are not always weather to be loved through; sometimes what presents as her mood is a claim: \emph{I need room; I need you to hear this; something between us is wrong.} A vow structured to love on regardless risks converting every such claim into one more contingency to be surmounted by steadfastness, so that the most resolute love becomes the love least able to \emph{hear}. The lover's constancy, a virtue from his side, can function from hers as a wall: her dissent is metabolized in advance as something his love is built to outlast. This is the gravest vector, for it disables the very channel through which a domination might be communicated and corrected, and it is the intimate face of the reasons-responsiveness hole of Section~\ref{sec:freewill}.

\textbf{(ii) Non-reciprocity.} Unconditional unilateral commitment distorts the economy of giving and receiving a just relationship maintains. A love that gives without condition and without reference to the other's giving can fix the beloved as the perpetual recipient of something she has not asked for in the form in which it arrives, or license the lover to endure, in fidelity's name, treatment he ought not to endure, unilaterally setting the terms of sacrifice as well. Reciprocity is not the demand that love be repaid like a debt; it is the condition that the shape of the giving be something both hands touch. Unilateral unconditionality removes her hand from it.

\textbf{(iii) Definitional capture.} Most subtly, the unilateral vow lodges the authority to define \emph{what loving her means} in the lover. As Section~\ref{sec:moral} showed, ``keep loving her'' requires content, and the vow assigns the filling of that content to one party. But what counts as loving this particular person is something over which she has standing to speak; to foreclose that, however lovingly, is to capture an authority that is properly shared, to love her in the third person while addressing her in the second.

None of these requires the lover to be selfish, deluded, or coercive. They are structural, following from the unilateral \emph{form} of the vow conjoined with the two-person character of what it governs. The vow can be free in the Kantian sense, sincere, constant, devoted, and still institute, through these channels, a quiet domination. The next section makes the structure precise by asking what kind of \emph{order} a vow institutes.

% =========================================================================
\section{The Vow as the Institution of a Rights-Order}\label{sec:rights}

To say a vow ``institutes an order'' can sound metaphorical. It is not. A vow over a shared life creates, between the two parties, a determinate structure of normative positions, and the most precise instrument for displaying that structure is Hohfeld's analysis of jural relations \parencite{hohfeld1919fundamental}. Hohfeld showed that what we loosely call ``a right'' decomposes into distinct positions that come in correlative pairs: a \emph{claim} in one party correlates with a \emph{duty} in the other; a \emph{privilege} (liberty) with a \emph{no-right}; and, at the second order, a \emph{power} to alter the parties' positions with a \emph{liability} to have one's positions altered, while an \emph{immunity} correlates with a \emph{disability}. Crucially, Hohfeld insisted that every such relation holds between exactly two persons: jural relations are irreducibly relational, never a property of one party alone.

\subsection{What the Unilateral Vow Distributes}

Read through this lens, the lover's vow is the unilateral institution of a set of Hohfeldian positions governing the relationship. The lover assumes \emph{duties} (to keep loving, whatever happens), and thereby confers on the beloved corresponding \emph{claims}. So far the vow looks generous: it loads obligation onto the lover and entitlement onto the beloved. But the decisive positions are at the second order. Who holds the \emph{power} to alter the terms, to determine what the duties and claims actually \emph{are}, what ``loving her'' shall be taken to require and provide? In the unilateral vow, that power is retained entirely by the lover. He defines the content; he may, being its author, redefine it; the beloved holds no correlative power but only a \emph{liability}, the position of one whose normative situation is determined by another's exercise of power. And does the beloved hold an \emph{immunity} against having the terms of the love set without her, correlating with a \emph{disability} in the lover to set them alone? She does not; the unilateral structure is precisely the lover's retention of the power and the beloved's exposure as liability.

This is the Hohfeldian translation of the non-domination diagnosis, and it sharpens it. The vow's generosity is all at the first order, claims and duties that favour the beloved, while the domination is all at the second order, in the distribution of power and immunity. One can be lavishly endowed with first-order claims and still be dominated, if one holds no power over the terms and no immunity against their unilateral revision. The benevolent master grants generous claims; what he withholds is the power and the immunity. Hohfeld lets us say exactly what a just intimate order would require that the unilateral vow omits: not more generous duties on the lover's side, but a \emph{distribution of the second-order positions}, a share of the power to set and revise the terms, and an immunity against their being set unilaterally. Justice in the rights-order is a matter of the second order.

\subsection{Order versus the Mere Imposition of Order}

The analysis also clarifies what is right about the impulse to vow. Some institution of an order is not optional in a shared life; two people who coordinate at all stand in jural relations, and the alternative to an order is not freedom but chaos, in which neither can rely on anything. The question is never whether to have a rights-order but \emph{how the power to constitute and revise it is distributed}. A vow that institutes an order while reserving all second-order power to its author has not done something gratuitous by instituting an order; it has done something unjust by \emph{monopolizing the constituent power}. This distinction, between the institution of an order and the monopolization of constituent power over it, is exactly the distinction we will need at the largest scale, and it is the bridge to the political-economic register: the same structure by which a benevolent legislator can institute an order that dominates those it generously binds recurs from the intimate dyad to the polity. Before drawing that isomorphism, we must ask how the legislator who monopolizes constituent power comes to do so without noticing, which returns us to the constitution of the legislating self.

% =========================================================================
\section{The Constitution of the Legislating Subject: Structuralist and Feminist Considerations}\label{sec:subject}

Every layer so far has deferred a question to this section. The Kantian noumenal refuge (Section~\ref{sec:kant}) presupposed a legislating self constituted prior to and independently of its relations. The reasons-responsiveness hole (Section~\ref{sec:freewill}) turned on a legislator already disposed, or not, to count the other as a source of reasons. The corrupted-inputs problem (Section~\ref{sec:justice}) and the monopolized constituent power (Section~\ref{sec:rights}) both pointed to a legislator whose sense of what love requires is given from somewhere. This section asks where, and the answer dismantles the assumption on which the purity of self-legislation depends. The argument runs as a chain: the legislating subject is not the origin but an effect (structuralism); the effect is produced by a determinate, historically gendered order (feminist political theory); therefore a sincere unilateral vow is liable to encode and reproduce that order's subordinations, which is why the three vectors of Section~\ref{sec:justice} are structurally expectable rather than personal lapses.

\subsection{The Subject as Effect, Not Origin}

The structuralist insight, in the form most useful here, is that the subject who appears to be the source of meaning and law is itself produced by structures that precede it. Althusser's account of \emph{interpellation} holds that individuals are constituted as subjects by being ``hailed'' by ideology: one becomes a subject by recognizing oneself in the positions a social order addresses to one, so that the very interiority from which one seems to legislate is an effect of one's having been addressed and having answered \parencite{althusser1971ideology}. Lacan's parallel claim is that the subject is an effect of its insertion into the symbolic order, the system of language and law that pre-exists it; the ``I'' that speaks and wills is constituted in and by a field of signifiers it did not author \parencite{lacan2006ecrits}. One need not adopt the whole of either system to take the structural point: the self that legislates the vow does not bring to the vow a conception of love minted by its own spontaneity; it brings a conception it has received, in which it has been formed as the kind of subject that loves in certain ways and expects certain things.

This is where the historicism of the reasons-responsiveness account (Section~\ref{sec:freewill}) returns with force. Fischer and Ravizza make the freedom-conferring character of a mechanism depend on \emph{how it was formed}: a mechanism instilled by manipulation does not ground freedom even if it now functions well \parencite{fischer1998responsibility}. The structuralist claim is that the mechanism of the loving self is \emph{always} formed by a process the self did not control, the process of subject-constitution itself. This does not by itself defeat freedom, on pain of making freedom impossible; but it means that whether the lover's legislating mechanism is genuinely his own, in the sense that grounds freedom and underwrites just authority, cannot be assumed and must be asked, and the asking requires examining the \emph{content} of the formation. What, concretely, is the historical structure that forms the subject who vows to love?

\subsection{The Structure Is Gendered Subordination}

For intimate partnership the answer is supplied by feminist political theory, and it is specific: the order that has historically constituted the loving subject is one of gendered subordination, and it has done so through the very institutions, marriage, the family, romantic love, that present themselves as the home of free and mutual affection. Pateman's analysis of the \emph{sexual contract} argues that the social-contract tradition's story of free and equal individuals contracting into civil society conceals a prior, sexual contract through which men secured a patriarchal right over women, and that the marriage contract in particular has been not a paradigm of free agreement between equals but a form of the subordination contract, in which one party contracts into a subordinate status under the language of consent \parencite{pateman1988sexual}. Okin's analysis of \emph{justice, gender, and the family} shows that the family, far from lying outside the scope of justice as a haven of spontaneous affection, has been a primary site of injustice, structured by a gendered division of labour and vulnerability that theories of justice, by exempting the family from scrutiny, helped to naturalize \parencite{okin1989justice}. And Young's account of the \emph{faces of oppression} insists that domination and oppression are not, in the first instance, matters of distributive shortfall or of individual ill will, but structural: they operate through the normal, often well-intentioned processes of ordinary life, including the everyday workings of intimate and familial roles \parencite{young1990justice}.

Put the structuralist mechanism and the feminist content together and the chain closes. The subject who legislates the vow has been constituted, as a loving subject, within a historically gendered order; the conception of love he brings to the vow, what love is, what it provides, what it may demand, what counts as the beloved's good, is not minted by his spontaneity but received from that order; and that order is one of subordination presenting itself as mutuality. Therefore the three vectors of Section~\ref{sec:justice} are not the idiosyncratic failings of a particular lover but the \emph{structurally expectable} output of a sincere subject legislating from a constituted standpoint he mistakes for a neutral one. Immunization against voice, non-reciprocity, and definitional capture are exactly what one would predict from a subject formed to experience a gendered allocation of emotional provision and authority as the natural shape of love. The lover need not be a patriarch in intention; he need only be a competent product of the order that formed him, legislating in good faith from inside it. This is why the noumenal refuge fails not merely as an epistemology but as a description: there is no constituted-independently legislator to retreat to, because the very capacity to legislate a conception of love is a deposit of the order under examination.

\subsection{Why This Does Not Collapse into Determinism}

It must be said at once what this argument does not claim, on pain of proving too much. It does not claim that the subject is wholly determined and incapable of legislating anything not dictated by the order that formed it; that would make critique, including this paper, impossible, and would dissolve the very agency whose freedom Section~\ref{sec:freewill} took seriously. The claim is weaker and more exact: there is no \emph{guaranteed-pure} standpoint, no noumenal refuge or non-relational practical identity, from which a self-sufficient and self-certifying legislation could issue. The capacity for critical distance is real but it is not the spontaneity of an uncaused legislator; it is itself a relational achievement, won through encounter with what does not fit the inherited conception, paradigmatically through the resistance of the other who says that what is offered as love is not, for her, love. This is the positive seed: if the formed subject cannot certify its own legislation from within, the correction can only come from without, from the standing of the other to contest, which Section~\ref{sec:regrounded} will make constitutive. First, though, we can now state precisely the line between a vow that institutes a legitimate order and one that institutes exploitation.

% =========================================================================
\section{Rights-Order or Exploitation? The Political Economy of the Vow}\label{sec:exploitation}

We can now answer the question that motivates the political-economic register: \emph{when} does the self-legislation of a shared life institute a legitimate rights-order, and \emph{when} does it institute exploitation? The materials are in hand: the second-order analysis of constituent power (Section~\ref{sec:rights}) and the constituted, gendered subject (Section~\ref{sec:subject}). What remains is to import, carefully, the formal structure of exploitation.

\subsection{Exploitation, Formally}

I take from the Marxian tradition the \emph{form} of exploitation while remaining agnostic about the labour theory of value, which the argument does not need. Formally, exploitation obtains when (a) one party systematically appropriates the benefit of another's contribution or labour, (b) the appropriating party, not the contributing one, controls the terms under which the contribution is rendered and counted, and (c) the arrangement is sustained by a representation that makes the appropriation appear natural, voluntary, or even loving, so that it is not experienced as appropriation at all \parencite[on the form of appropriation and its mystification, see][esp. the analysis of the wage form]{marx1867capital}. Condition (c), the ideological masking, is not incidental; it is what distinguishes exploitation from open coercion and what makes it stable. Marx's point about the wage form, that it presents as an equal exchange (a fair day's wage for a fair day's work) a relation in which surplus is appropriated, has its exact analogue in the intimate sphere, where the language of love presents as the free overflow of affection a relation in which one party's relational labour is appropriated without recognition or reciprocal power.

\subsection{Love as the Site and the Mask}

The feminist development of this analysis, in social-reproduction theory, supplies the intimate content. The work of sustaining persons emotionally and materially, care, attention, the management of feeling, the labour of keeping a relationship and a household going, is \emph{labour}, even though it is performed outside the wage and is therefore, in the theory's terms, not \emph{exploited} in the strict wage sense but \emph{expropriated}: necessary work that is unpaid, undervalued, and rendered invisible by being recoded as the natural expression of love or femininity \parencite{federici2012revolution, fraser2016contradictions}. Federici's central claim is exactly that housework and care were made to appear as expressions of love rather than as work, and that this recoding is the mechanism of their expropriation \parencite{federici2012revolution}. Hochschild's analysis of \emph{emotional labour} names the further dimension: the management of one's own feeling to produce a required emotional state in another is work, and it is distributed unequally, with one party characteristically assigned the open-ended task of sustaining the relationship's emotional weather \parencite{hochschild1983managed}. Social-reproduction theory's distinction between exploitation and expropriation, and its insistence that the sphere of love and care is where capitalism offloads the costs of reproducing people, give the intimate analysis its political-economic depth \parencite{bhattacharya2017social}.

Now the unilateral vow can be located precisely on this map. A vow that reserves second-order power to the lover, made by a subject constituted within a gendered order, and expressed in the language of unconditional love, satisfies all three conditions of the formal structure of exploitation when it operates so as to enclose the beloved as the unrecognized provider of relational and reproductive labour. (a) The benefit of her sustaining labour is appropriated, the relationship is kept going, the lover is emotionally and materially sustained, while (b) the terms, what counts as love, what is owed, what may be asked, are controlled by him through his retained constituent power, and (c) the whole is masked as love, indeed as \emph{his} love for \emph{her}, so that her labour is recoded as the natural expression of her affection and disappears as labour. The very unconditionality that made the vow look like the height of devotion is, on this analysis, the perfect ideological form: it presents as pure gift a structure of appropriation, and it does so most effectively precisely when the lover is sincere, because his sincerity is the surface under which the structure operates unseen.

\subsection{The Line}

This yields the determinate line the section promised. The self-legislation of a shared life institutes a \emph{legitimate rights-order} when the other retains co-authorship over the second-order terms, the power to set and revise what is given and demanded, and an immunity against having those terms fixed unilaterally; under those conditions the labour of sustaining the relationship is rendered and counted on terms the contributor co-controls, and is recognized as contribution. The same self-legislation institutes \emph{exploitation} when the other is enclosed as the provider of unrecognized relational labour under terms she does not co-control, with the enclosure masked by the language of love. The line is not drawn by the warmth of the feeling, the generosity of the first-order provisions, or the sincerity of the legislator, all of which are compatible with exploitation and indeed are its most effective cover. It is drawn, exactly as in the political-economic case, by the distribution of constituent power and by whether the contributor's labour is recognized as such or recoded as natural devotion. The intimate and the political-economic are here not analogous but, in the relevant structural respect, identical: in both, the difference between a just order and an exploitative one is whether those whose labour sustains the order retain the power to set its terms and the standing to have their contribution recognized.

% =========================================================================
\section{The Objection from Ideological Recognition}\label{sec:ideology}

The argument now points unmistakably toward a remedy: make the vow relational; distribute the second-order power; give the beloved constitutive standing to contest and co-author what loving her means. But the most serious objection to that remedy comes from within the very traditions, recognition theory and relational autonomy, on which it draws, and it must be met before the remedy can be stated, on pain of the remedy's being hollow.

The objection is that relationality is no guarantee of justice, because recognition can itself be \emph{ideological}. Honneth, whose theory makes love the first sphere of recognition, the sphere in which basic self-confidence is formed, and recognition the medium of self-realization \parencite{honneth1995struggle}, is himself the source of the difficulty. He observes that subjects can derive genuine self-worth and stable identity from recognition conferred within, and securing their consent to, a subordinate position \parencite{honneth2007disrespect}. The devoted partner may attain a robust sense of identity and esteem precisely by fulfilling a role that subordinates her; her recognition is real and her endorsement sincere, yet the arrangement she endorses is one we have reason to call unjust. ``Ideological recognition'' names the cases where being recognized, and welcoming it, coexists with and even cements domination. Pressed against relational autonomy, the point becomes a dilemma: \textcite{khader2020feminist} argues that the social embedding relational theories celebrate is also what can produce a person's sincere attachment to conditions that diminish her, the problem of \emph{adaptive preferences} \parencite[cf.][]{khader2011adaptive}. Applied to our case: the lover opens the terms to co-authorship, and the beloved, formed by the very gendered order Section~\ref{sec:subject} described to ask for little, co-authors terms that subordinate her and cherishes them. The relationship is now relational, recognitional, co-authored, and unjust, the more insidiously because it wears the appearance of mutual endorsement.

This is a genuine knife at the throat of the remedy, and it is sharpened by the present paper's own Section~\ref{sec:subject}: if the subject is constituted by a gendered order, why expect the \emph{beloved's} constituted preferences to be any purer a touchstone than the lover's constituted conception? I do not blunt the objection by denying its force. Turning the vow relational does not, by itself, secure justice. If the criterion of a just relational order were \emph{de facto endorsement}, that both parties in fact accept the terms, ideological recognition would defeat the remedy outright, since de facto endorsement is exactly what the contented subordinate supplies.

\subsection{Contestability, Not Endorsement}

The answer is that the operative criterion is not de facto endorsement but \emph{live contestability}, and this is where the republican analysis of Section~\ref{sec:justice} pays off. Recall that domination is modal: it concerns not whether interference \emph{actually} occurs but whether one is exposed to its possibility at another's discretion \parencite{pettit1997republicanism}. Transpose this to recognition. What makes a relational order just is not that its terms are in fact accepted, but that each party retains a modally robust standing to contest and revise them, a standing tested not by whether dissent actually occurs but by what \emph{would} happen if it did. The contented subordinate's situation is unjust, on this criterion, not because she fails to object, she does not wish to, but because, \emph{were} she to object, her objection would carry no weight; the order is structured so that her dissent would not move it. Endorsement under conditions where dissent is foreclosed is not the kind of endorsement that confers justice. Contestability asks the counterfactual, modal question, is dissent live, weighty, efficacious across the relevant range of circumstances?, and that question discriminates between recognition that secures standing and recognition that merely secures consent to its absence. This is precisely why de facto endorsement fails as a criterion and modal contestability succeeds: the former is a fact about the actual world that ideology can manufacture, the latter a fact about the structure of power across possible worlds that ideology cannot, by manufacturing actual consent, satisfy.

Two clarifications fortify the reply. First, contestability is not satisfied by a merely formal permission to object. A relationship in which the beloved is ``allowed'' to disagree but in which her disagreement is structurally inert, always outlasted, always reframed, never able to change anything, fails the test as surely as one that forbids objection outright; this is the relational analogue of Hohfeld's point that a nominal claim without a correlative power is no protection. Live contestability requires that dissent be \emph{able to make a difference}, a fact about the distribution of second-order power and of responsiveness, not about the presence of a nominal right. Second, following Ricoeur, contestability presupposes a \emph{proper distance} between the parties: a relation close enough for intimacy yet preserving the irreducible separateness of the other, so that the beloved remains an other who can say no rather than being absorbed into the lover's conception of the shared good \parencite{ricoeur2005course}. Recognition collapses into ideology precisely when this distance collapses, when the other is so identified with the relationship, or with the lover's image of her, that there is no standpoint from which she could dissent. Proper distance is the structural condition that keeps contestability live, and it connects back to the constitution of the subject: the maintenance of a standpoint from which the other can contest is what prevents her constituted preferences from being simply the internalized echo of the lover's, or of the order's.

\subsection{The Honest Residue}

In keeping with the standard of not over-claiming, the limits of the reply must be stated. Contestability is not an algorithm and not a guarantee. There are hard cases in which it is genuinely unclear whether dissent is live or foreclosed, and cases in which a deeply internalized subordination corrodes the very capacity to conceive of dissent, so that the counterfactual test threatens to return indeterminate, the standpoint of contestation having been, as it were, colonized in advance. I do not claim that contestability dissolves the problem of ideological recognition. I claim that it locates justice in the right place, in the modally robust standing to contest rather than in the fact of consent, and that it does decisively better than de facto endorsement, which ideology defeats immediately. Where the capacity to contest has itself been damaged, the requirement of \emph{proper distance} and the relational achievement of critical standpoint (Section~\ref{sec:subject}) become the site of repair, slow, uncertain, and never to be presumed complete. A criterion can be correct and still leave a residue of hard cases; epistemic honesty requires naming the residue rather than papering over it.

% =========================================================================
\section{The Vow Re-Grounded: Relational Co-Legislation and Generative Justice}\label{sec:regrounded}

We can now reconstruct the vow so as to preserve the goods established along the way, the freedom of Section~\ref{sec:freewill}, the moral worth of Section~\ref{sec:moral}, while closing the gap diagnosed across Sections~\ref{sec:justice}--\ref{sec:exploitation} and surviving the objection of Section~\ref{sec:ideology}. The reconstruction does not repudiate the original vow. Its core, that love should not be the hostage of mood or the dependent variable of reciprocation, is worth keeping; it is what made the vow an achievement of freedom. What must change is not the constancy but the unilaterality, not the first-order devotion but the monopoly of second-order power.

The original vow was: \emph{Whatever happens, however she feels, I will keep loving her.} The reconstructed vow is closer to:

\begin{quote}\itshape
I will keep loving her, and what loving her means, here and now, I hold open to her: to her voice, her dissent, her revision. My constancy is that I will not make my love hostage to her mood or contingent on her return; but my constancy does not stand above her standing to contest what my love does to her and asks of her. The terms of our love are hers to set with me, and never mine to set alone.
\end{quote}

The first clauses preserve the freedom layer: the love remains self-legislated, withdrawn from the jurisdiction of mood and reciprocity. The reasons-responsiveness criterion of Section~\ref{sec:freewill} is now satisfied rather than violated, since the mechanism is explicitly held open to the beloved as a source of reasons. The final clauses close the justice gap by relocating the \emph{second-order power} (Section~\ref{sec:rights}) from the lover alone to the two together: constancy is retained as a property of the lover's commitment to remain and to keep willing her good, while what is surrendered is the monopoly over the interpretation of that good. The three vectors are addressed at their root: voice is constitutively admitted, not immunized against; the shape of giving becomes something both hands touch; the definition of loving her is restored to shared authority. And because the operative criterion is live contestability rather than endorsement, the reconstruction does not collapse into ideological recognition: the test is not whether she accepts the terms but whether her dissent remains live, weighty, and able to move them, across the proper distance that keeps her an other who can say no. The political-economic line of Section~\ref{sec:exploitation} is thereby crossed back to the right side: with second-order power shared and her sustaining labour recognized as contribution rather than recoded as natural devotion, the order ceases to be exploitative and becomes a genuine rights-order.

\subsection{From Self-Legislation to Co-Legislation}

Mark what has happened to the concept of self-legislation. We began with a self-standing rational agent giving himself a law. We end with a law that governs the lover's love and is not, and cannot be, his alone, a law that is, at its base, \emph{co}-legislation. This is not a concession forced from outside onto an otherwise solitary act; Section~\ref{sec:subject} showed that the solitary act was never solitary, that the legislating self was constituted relationally all along. The reconstruction does not \emph{add} relationality to a self-sufficient legislation; it brings the legislation's own grounds into view and aligns the form of the vow with the truth of its origin. The Kantian picture presupposed a legislator constituted prior to and independently of relations; the relational ontology that this paper has progressively established, through the failure of the noumenal refuge, the historicism of agency, and the structural constitution of the subject, holds instead that the loving subject is in part \emph{generated in the relation}: his sense of what love is, what the other is owed, what he himself wants and is, takes shape in and through the relationship he proposes to legislate over \parencite{nedelsky1989reconceiving, mackenzie2000relational, helm2010love}. ``Self''-legislation in the intimate sphere was always a moment within the ongoing co-constitution of two subjects by each other.

\subsection{Generative Justice}

This licenses, at the end and not before, a conception of justice thicker than the non-domination floor that did the diagnostic work. If the subjects are generated in the relation, then justice between them is not merely the absence of domination but the positive maintenance of conditions under which each can continue to become who they are in the relation without that becoming being captured or foreclosed by the other. Call this \emph{generative justice}: justice measured by whether the relation keeps open the conditions of each party's ongoing self-formation, including, centrally, the standing to contest the terms of the relation itself. Non-domination and live contestability are the floor; generative justice is the substantive ideal they point toward, and it is the proper aim of the reconstructed vow. The reason to introduce it only here, rather than to build the whole argument on it, is methodological honesty: the diagnosis and the reply stand on the thin, widely shared notion of non-domination alone, and do not require the reader to accept the thicker, relational-ontological conception. Generative justice names what the reconstructed vow is \emph{for}; non-domination and contestability are what it must, at minimum, secure.

% =========================================================================
\section{Conclusion: The Isomorphism of the Vow and the Order}\label{sec:conclusion}

This paper began with a lover's vow and has arrived at a structure that the vow shares with orders far larger than itself. The vow, taken as an instance of the self-legislation of a shared life, was shown to be an exemplary exercise of Kantian freedom whose freedom, on any operable criterion, is conditional on its remaining responsive to the other; a morally admirable form whose content is silent and capturable; and, in its unilateral version, an institution of domination in the modal sense, a benevolent-master order generous at the first order and dominating at the second. A Hohfeldian analysis located the wrong precisely in the monopolization of second-order constituent power; an account of the constituted subject showed why a sincere legislator reproduces that wrong without intending it; and the political economy of social reproduction drew the line between the vow that institutes a rights-order and the vow that institutes exploitation, the appropriation of unrecognized relational labour under the mask of love. The remedy is not a purer legislator but a relational co-legislation whose justice rests on the other's modally robust, live standing to contest, and whose aim is the generative justice that keeps open each party's becoming.

The deepest yield is the isomorphism the introduction promised. The micro-order a lover institutes by a vow and the macro-order a political economy institutes by its distribution of rights and labour are, in the structural respects that matter, the same. In both, a legislation that presents itself as self-grounded is in fact grounded in relations and a history that precede it. In both, procedural correctness on the legislating side, sincerity, universalizability, good will, generous first-order provision, cannot certify the justice of the order, because justice is a matter of the second order, of who holds the power to set and revise the terms and whether those bound retain the standing to contest them. In both, the difference between a legitimate order and domination or exploitation is not the benevolence of the powerful but the secured, non-discretionary standing of those whose participation and labour sustain the order. And in both, the ideology that most effectively masks domination is the one that recodes appropriation as gift, whether as the wage form's ``fair exchange'' or as love's ``unconditional devotion.''

That a lover's most earnest vow and a political economy's distribution of rights should share a structure is not a deflation of love to economics, nor an inflation of economics to romance. It is a claim about normativity as such: that the standing of any law over a shared life, from the smallest dyad to the largest polity, rests not on the purity or the goodwill of its author but on the irreducible standing of those it binds to answer it. The lover who understands this does not love less freely or less faithfully. He loves, for the first time, justly, having understood that the law of a shared life was never his to give alone, and that the deepest constancy is not the will that holds its course against the beloved, but the will that keeps, between them, the room in which she may always answer.

\medskip
{\color{warnred}\itshape What is built between two people, an order, a vow, a shared life, is never neutral and never merely given: it either secures for both the standing to remake its terms, or it quietly fixes those terms as one party's gift and the other's fate. To keep loving, then, is not to hold the order steady against her, but to keep handing back to her, again and again, the power to make it ours.}

% =========================================================================
\section*{Acknowledgements}
\addcontentsline{toc}{section}{Acknowledgements}

This paper is the third in a series on the philosophy of intimacy and the theory of justice. The lover and beloved discussed here are kept deliberately anonymous, and the case is presented as a constructed one, so that the argument may stand on its structure rather than on any private history.

My deepest gratitude is owed to the one I have loved all along, the ``forest girl'' to whom the first paper of this series was already devoted: pure and natural in spirit, good-hearted, resilient, and wise, carrying within her a quiet sense of mission toward the world. She is the motivation for the writing of this paper, as she has been for so much else. The vow examined in these pages was occasioned by what she awakened; and the paper's conclusion, that the deepest constancy is not the will that holds its course against the beloved, but the will that keeps open the room in which she may always answer, is the form in which I renew that vow to her: that I will not fail her, to the end of our days.

In the interest of transparency, I note that an AI assistant was used in preparing this manuscript, as a tool for drafting, structuring, and refining the argument and its prose; the ideas, commitments, and final judgements are my own, and I take full responsibility for the content.

\begin{center}
{\cjkfont\color{rosedeep} 愿天下有情人，皆得相敬如宾，权责相与，白首不相负。}
\end{center}

% =========================================================================
\printbibliography[heading=bibintoc]

\end{document}
